% =====================================================================
\chapter{System Analysis and Requirements}
\label{ch:analysis}

\section{Stakeholder Analysis}
Four groups of stakeholders interact with the system, each with distinct
needs:

\begin{description}
  \item[Passengers (students and staff).] The primary beneficiaries. They
  need to see the live position of their bus, the estimated arrival time
  at \emph{their own} stop, the day's schedules, and timely alerts---in
  English or Bangla according to their preference. They interact only
  with the mobile application.
  \item[Drivers.] Deliberately, drivers have \emph{no required role} in
  normal operation: the fixed tracker reports automatically. A driver
  participates only in the optional fallback path, by carrying a
  smartphone that publishes location when a bus has no working tracker.
  \item[Transport administrators.] They manage routes, stops, schedules,
  buses, GPS devices, and driver assignments, and monitor the whole fleet
  on a live operations map. They interact with the web administration
  console.
  \item[The institution.] The university requires the recurring cost to
  stay within a modest budget, and the deployment to run on its real
  fleet without disrupting operations.
\end{description}

\section{Functional Requirements}
The functional requirements were derived from the stakeholder analysis
and refined during early development. Table~\ref{tab:freq} lists the
principal requirements; each is referenced from the design chapter where
it is satisfied.

\begin{table}[H]
\caption{Principal functional requirements.}
\label{tab:freq}
\centering
\small
\begin{tabular}{@{}llp{9.2cm}@{}}
\toprule
\textbf{ID} & \textbf{Area} & \textbf{Requirement} \\
\midrule
FR-1 & Acquisition & Ingest positions automatically from a fixed
vehicle GPS tracker with no driver involvement. \\
FR-2 & Acquisition & Accept authenticated position reports from a driver
smartphone as an optional secondary source. \\
FR-3 & Acquisition & Select the best available source per fix,
automatically and without passenger-visible artefacts. \\
FR-4 & Quality & Reject inaccurate, physically implausible, and
jittery fixes before they affect the public position. \\
FR-5 & ETA & Compute an arrival estimate for every upcoming stop of the
active route on each accepted fix, with a confidence label. \\
FR-6 & Delivery & Push position and ETA updates to all subscribed
passengers within seconds over a persistent connection. \\
FR-7 & Lifecycle & Show the bus state before departure (parked,
approaching origin, arrived at origin) instead of a blank map. \\
FR-8 & Lifecycle & Mark displayed data as stale when the latest fix ages
beyond a threshold. \\
FR-9 & Alerts & Notify a subscribed rider when the bus reaches an 80\,m
radius of a stop, including their chosen drop-off stop. \\
FR-10 & Localization & Provide the complete passenger experience in both
English and Bangla, switchable in-app. \\
FR-11 & Administration & Manage routes, stops, schedules, buses, GPS
devices, and user roles through a web console. \\
FR-12 & Administration & Display all active vehicles on a live
operations map for transport staff. \\
FR-13 & Audit & Log every raw report, every accepted fix, and every
ETA prediction for later analysis. \\
\bottomrule
\end{tabular}
\end{table}

\section{Non-Functional Requirements}
\begin{table}[H]
\caption{Principal non-functional requirements.}
\label{tab:nfreq}
\centering
\small
\begin{tabular}{@{}llp{9.2cm}@{}}
\toprule
\textbf{ID} & \textbf{Quality} & \textbf{Requirement} \\
\midrule
NFR-1 & Latency & A position fix should reach subscribed passengers
within seconds of capture under normal network conditions. \\
NFR-2 & Resilience & Loss of any single position source must degrade
coverage gracefully, not eliminate it. \\
NFR-3 & Cost & The entire backend must run on one low-cost virtual
private server; no managed cloud services are required. \\
NFR-4 & Efficiency & Reporting and polling cadence must not waste
cellular data or vehicle battery while buses are parked. \\
NFR-5 & Security & All client--server communication is authenticated;
role-based authorization separates passengers from administrators. \\
NFR-6 & Scalability & Realtime fan-out must be able to scale
horizontally across server processes without losing consistency. \\
NFR-7 & Usability & The passenger app must remain informative during
edge conditions (pre-trip, signal loss, stationary bus). \\
\bottomrule
\end{tabular}
\end{table}

NFR-4 deserves a remark: it was \emph{added during the pilot}, when
operational data revealed that roughly 98\% of raw device reports were
being transmitted while buses were parked. The discovery and the
resulting optimization are treated fully in
Section~\ref{sec:eval-efficiency}; the requirement is recorded here
because it reshaped both the device configuration and the server's
polling design.

\section{Feasibility Study}

\subsection{Technical Feasibility}
Every component of the proposed system relies on proven, freely available
technology: a mass-produced vehicle GPS/GSM tracker (Concox GT06S), an
open-source telematics server (Traccar) that already implements the
tracker's binary protocol, an open-source relational database
(PostgreSQL), an in-memory store (Redis), and mature application
frameworks (Node.js/Express for the server, React Native/Expo for the
mobile app, Next.js for the console). No component required custom
hardware or proprietary licensing, so technical feasibility was
established by construction.

\subsection{Economic Feasibility}
The recurring infrastructure is a single low-cost virtual private server
hosting the application server, Traccar, PostgreSQL, and Redis together.
The one-time hardware cost is the tracker unit and its SIM per bus. Both
figures sit comfortably within a student-project budget, and---more
importantly for adoption---within what a university transport office
could sustain for a full fleet. The economic design goal, ``everything on
one server, no external paid services,'' held throughout the pilot.

\subsection{Operational Feasibility}
Operational feasibility hinged on one question: can the system run
without adding work for anyone? The hardware-first design answers it. The
tracker is hard-wired to the bus and reports automatically; drivers need
not carry, charge, or operate anything; administrators interact with the
console only to maintain routes and schedules, which they already did on
paper. The seven-week unattended pilot (Chapter~\ref{ch:evaluation})
confirmed this in practice.

\section{Development Methodology}
The system was developed iteratively. An initial vertical slice---one
bus, one route, raw positions on a map---was deployed early, and each
subsequent iteration added one capability (quality filtering, ETAs, trip
lifecycle, notifications, bilingual UI, the administration console)
followed by observation on the live fleet. Two practices shaped the
result. First, \emph{operational data drove priorities}: the staleness
indicator, the pre-trip map state, and ultimately the idle-reporting
optimization all originated from watching real deployment behaviour
rather than from the original plan. Second, \emph{every layer logs}: raw
ingests, accepted fixes, trip events, and ETA predictions are all
persisted, which is what later made the evaluation in
Chapter~\ref{ch:evaluation} possible directly from the production
database.

\section{Tools and Technologies}
Table~\ref{tab:stack} summarizes the technology choices; the rationale
for each appears in the corresponding design section.

\begin{table}[H]
\caption{Technology stack of \system{}.}
\label{tab:stack}
\centering
\small
\begin{tabular}{@{}lll@{}}
\toprule
\textbf{Layer} & \textbf{Technology} & \textbf{Role} \\
\midrule
Vehicle & Concox GT06S & Fixed GPS/GSM tracker (primary source) \\
Vehicle & Android smartphone & Optional driver-side fallback source \\
Ingestion & Traccar & Tracker protocol decoding, device management \\
Server & Node.js + Express (TypeScript) & Application server \\
ORM & Prisma & Typed access to the relational schema \\
Database & PostgreSQL & System of record (31 models) \\
Cache/bus & Redis & Hot-path cache and pub/sub back-plane \\
Realtime & Socket.IO & WebSocket rooms per active trip \\
Mobile & React Native + Expo & Bilingual passenger application \\
Mobile & React Native + Expo & Driver app (fallback publisher) \\
Web & Next.js (React) & Administration console \\
Maps & Google Maps APIs & Map rendering and route building \\
Hosting & Single VPS & All backend components self-hosted \\
\bottomrule
\end{tabular}
\end{table}

\section{Summary}
This chapter identified the system's four stakeholder groups, fixed
thirteen functional and seven non-functional requirements, established
technical, economic, and operational feasibility, and recorded the
iterative, data-driven development methodology and technology stack. The
next chapter presents the design that satisfies these requirements.
